\section{Modular inverse}

The modular inverse of $a$ modulo $m$ is a value $x$ satisfying $ax \equiv 1 \pmod m$. It exists exactly when $\gcd(a,m)=1$. Modular division means multiplication by an inverse: $a/b \pmod m = a \cdot b^{-1} \pmod m$, only when $b$ is invertible.

\subsection{Prime modulus: Fermat's little theorem}
For prime \texttt{MOD} and \texttt{a \% MOD != 0}, $a^{-1} \equiv a^{\texttt{MOD}-2} \pmod{\texttt{MOD}}$.

\begin{lstlisting}
using ll = long long;
const ll MOD = 1000000007LL;

ll modPow(ll a, ll e) {
    ll ans = 1; a %= MOD;
    while (e) {
        if (e & 1) ans = ans * a % MOD;
        a = a * a % MOD;
        e >>= 1;
    }
    return ans;
}
ll modInverse(ll a) { // MOD must be prime; a must not be 0 mod MOD
    a %= MOD;
    if (a < 0) a += MOD;
    return modPow(a, MOD - 2);
}
\end{lstlisting}

\subsection{Arbitrary modulus: extended Euclid}
For a modulus that may be composite, solve $ax+my=1$. The resulting $x$ is the inverse if and only if the gcd is 1.

\begin{lstlisting}
ll extendedGCD(ll a, ll b, ll& x, ll& y) {
    if (!b) return x = 1, y = 0, a;
    ll x1, y1, g = extendedGCD(b, a % b, x1, y1);
    x = y1; y = x1 - (a / b) * y1;
    return g;
}
ll inverseAnyMod(ll a, ll mod) {
    ll x, y;
    if (extendedGCD(a, mod, x, y) != 1) return -1;
    return (x % mod + mod) % mod;
}
\end{lstlisting}

\section{Working under \texttt{1e9 + 7}}

\texttt{1000000007} is prime, so Fermat inverses work for non-zero residues. Use \texttt{long long} for products: two residues fit safely in 64 bits. For moduli near $10^{18}$, multiply with \texttt{\_\_int128}.

\begin{lstlisting}
ll addMod(ll a, ll b) { a += b; return a >= MOD ? a - MOD : a; }
ll subMod(ll a, ll b) { a -= b; return a < 0 ? a + MOD : a; }
ll mulMod(ll a, ll b) { return a * b % MOD; }
ll mulBigMod(ll a, ll b, ll mod) { return (__int128)a * b % mod; }
ll negMod(ll a) { a %= MOD; if (a < 0) a += MOD; return (MOD - a) % MOD; }
\end{lstlisting}

Never write \texttt{(a / b) \% MOD}; use \texttt{a * modInverse(b) \% MOD} after checking \texttt{b \% MOD != 0}. Normalize negative results in C++. You may cancel a factor modulo $m$ only when that factor is invertible. A matrix inverse and a function inverse are different concepts from a multiplicative modular inverse.

\section{Efficient matrix multiplication}

For $A$ of size $n \times m$ and $B$ of size $m \times p$, their product has $C_{ij}=\sum_k A_{ik}B_{kj}$. The cache-friendly \texttt{i-k-j} loop reuses \texttt{A[i][k]} and can skip zero entries.

\begin{lstlisting}
using Matrix = vector<vector<ll>>;
Matrix multiply(const Matrix& A, const Matrix& B) {
    int n = A.size(), m = B.size(), p = B[0].size();
    Matrix C(n, vector<ll>(p));
    for (int i = 0; i < n; i++) for (int k = 0; k < m; k++) {
        ll aik = A[i][k];
        if (aik == 0) continue;
        for (int j = 0; j < p; j++)
            C[i][j] = (C[i][j] + aik * B[k][j]) % MOD;
    }
    return C;
}
\end{lstlisting}

Multiplication costs $O(N^3)$ for square $N \times N$ matrices. It is practical in matrix exponentiation when the state dimension is small (often 2--50), including linear recurrences, fixed-transition DP, and walks of exactly $K$ graph edges. For fixed dimensions, \texttt{array} can be faster than dynamic vectors.

\section{Factorials, combinations, and permutations}

For many combinatorics queries under a prime modulus, precompute factorials and inverse factorials. This requires the maximum queried $n$ to be strictly smaller than \texttt{MOD}.

\begin{lstlisting}
const int MAXN = 1000000;
vector<ll> fact(MAXN + 1), invFact(MAXN + 1);

void initCombinatorics() {
    fact[0] = 1;                         // 0! = 1
    for (int i = 1; i <= MAXN; i++) fact[i] = fact[i - 1] * i % MOD;
    invFact[MAXN] = modPow(fact[MAXN], MOD - 2);
    for (int i = MAXN; i >= 1; i--) invFact[i - 1] = invFact[i] * i % MOD;
}
ll nCr(int n, int r) {
    if (r < 0 || r > n) return 0;
    return fact[n] * invFact[r] % MOD * invFact[n - r] % MOD;
}
ll nPr(int n, int r) {
    if (r < 0 || r > n) return 0;
    return fact[n] * invFact[n - r] % MOD;
}
\end{lstlisting}

Precomputation is $O(N)$, then \texttt{nCr} and \texttt{nPr} are $O(1)$. Use them for choosing or arranging objects, binomial coefficients, grid paths, inclusion-exclusion, and counting formulas. If $n \geq p$ under prime modulus $p$, then $n! \equiv 0 \pmod p$, so the usual factorial-inverse formula fails; check Lucas' theorem or another advanced method.

\subsection{All inverses in linear time}
For prime \texttt{MOD} and $1 \leq i < \texttt{MOD}$:
\begin{lstlisting}
vector<ll> inv(N + 1); inv[1] = 1;
for (int i = 2; i <= N; i++)
    inv[i] = MOD - (MOD / i) * inv[MOD % i] % MOD;
\end{lstlisting}

\section{Recognition and checklist}
\begin{itemize}
\item Division, probabilities, expected values, \texttt{nCr}, or fraction-like expressions modulo a prime: \textbf{modular inverse}.
\item Inverse modulo an arbitrary/composite modulus: \textbf{extended Euclid}, after checking gcd 1.
\item Huge linear recurrence or a small fixed state transition: \textbf{matrix multiplication + exponentiation}.
\item Many \texttt{nCr}/\texttt{nPr} queries: \textbf{factorial + inverse factorial}.
\item $n$ or $r \geq \texttt{MOD}$: do not blindly use factorial inverses; consider \textbf{Lucas' theorem}.
\end{itemize}

Frequent errors: dividing directly modulo \texttt{MOD}; forgetting to normalize subtraction; asking for inverse of zero; applying Fermat with a composite modulus; using \texttt{int} for products; forgetting \texttt{0! = 1}; and allocating factorial arrays beyond the needed maximum.
